\section{Evaluation Method}
\label{sec:evaluation}

\subsection{Scope and design}

We conduct a software-system evaluation of the contract rather than a study of
generated language or human experience.  The design has four linked parts:
(1) a migration and coverage audit over every modeled scene asset, (2) a
controlled non-regression comparison with the direct artifacts from which the
contract instances were generated, (3) executable backend and Unity contract
tests, and (4) mutation and timing analyses that separate common faults from
contract-only cross-layer faults.  All reported proportions use exact
numerators and denominators for the fixed corpus.

The corpus contains the three spatial development cases available in the
repository: a career-fair room, a classroom with a dinosaur exhibit, and a
food/trade-fair brand room.  They contain 27, 36, and 30 routable scene assets
and 5, 4, and 6 modeled agents, respectively.  The cases were authored during
system development, not sampled from a population or held out from
implementation decisions.  This scope supports exhaustive checks over the
three models, not generalization to other rooms or authors.

\subsection{RQ1: coverage and non-regression}

For each case, the benchmark regenerates a fresh native V2 instance in memory
from its frozen version-0.5 case model.  Checked-in V2 files are hashed and
audited but are not used as treatment input.  An asset is complete only if it
appears in at least one asset-targeting interaction chain and every such chain
resolves exactly one capability and provider, agrees with the scenario
performer, contains the evaluated asset among its targets, reaches an
executable runtime binding and action, resolves all resulting assertions, and
has a use-case-linked validation case covering the binding and assertions.
The asset denominator contains every V2 entity with
\texttt{entityRole=sceneObject}; the chain denominator contains every
\texttt{ScenarioStep}/\texttt{CapabilityUse} pair explicitly targeting one of
those assets.  A complete chain is a declaration and traceability result, not
evidence that a conversation was executed.

The non-regression comparison uses two independently parsed adapters over the
same frozen cases:

\begin{description}
  \item[Direct wiring.] Reads the existing scene-semantics, generated
  agent-role, and handoff-matrix artifacts.
  \item[Contract representation.] Reads the fresh V2 instance regenerated in
  the same benchmark run.
\end{description}

Both expose a normalized ownership graph and the same decision policy:
asset assignment outranks object-group assignment, which outranks zone
responsibility; more than one owner at the highest active tier is ambiguous.
For every routeable object, the benchmark compares the owner and then classifies
a route from every modeled starting agent as local, directly permitted,
transitively graph-reachable, or unreachable.

The parsers extract candidate sets and handoff edges separately, but one shared
routing function classifies both normalized views.  Agreement is consequently
a migration and serialization non-regression check under a fixed policy, not
independent evidence that the policy itself is correct.

All natural objects currently have asset-level owners.  To exercise the entire
precedence rule without misrepresenting corpus frequencies, a separate derived
suite copies one object per case and creates three probes: remove its asset
owner to require a group fallback; remove both asset and group owners to
require a zone fallback; and add a competing owner at the highest active tier
to require ambiguity rejection.  The direct artifacts have no native group
relation, so the benchmark inserts the same normalized group candidate into
both copies.  These probes test decision behavior, not equivalent native
authoring expressiveness.

\subsection{RQ2: fail-closed runtime enforcement}

The exhaustive runtime sweep converts all object-by-start-agent combinations
into backend requests, for 459 probes in total.  The expected owner and route
class come from the frozen direct-wiring projection.  A fresh V2 project is
materialized and verified equal to the in-memory treatment before execution.
Each probe runs through \texttt{SessionStore.chat} with a deterministic
structured-response stub and an active network guard.

Local and direct routes are expected to be admitted.  An admitted response must
preserve the expected target and provider and contain non-empty,
cross-field-consistent capability-use, capability, binding, and action
identifiers.  Transitive and unreachable routes are expected to be rejected
because one request permits at most one direct handoff.  A rejected request
must stop before the stub and leave session state unchanged.  The benchmark
restores a session snapshot after every probe so that cases remain independent.

A targeted spatial-contract suite exercises one valid deictic request and 16
invalid variants, including missing context, unknown or contradictory
identifiers, obsolete hashes, ambiguous candidates, invalid active agents,
oversized fields, non-finite coordinates, and an undeclared handoff proposed by
the stub.  Request-decidable variants require zero stub calls and no retained
mutation.  The undeclared handoff is knowable only after generation; it
requires one stub call and must then roll back provisional state.

Three Unity Editor batch entry points load the same generated artifacts and
runtime components used by the client.  They exercise scenario progression,
dispatch, five assertion types, valid and broken evidence, desktop-ray
selection, ambiguous binding registries, exact target-specific multi-scenario
selection, and permission checks.  Exit status and explicit success markers are
recorded.  These are executable Editor paths, not participant or headset
trials.

\subsection{RQ3: fault scope, representation edits, and local cost}

The \emph{common} mutation suite applies three changes that both
representations can express, once per case: remove all ownership paths for one
object, add a second asset-level owner, and replace a handoff target with a
dangling identifier.  Each mutation starts from an expected-valid copy.
A detection counts only a new issue relative to that copy; localization
requires the new issue to name a mutation-specific object, agent, or field
token.  We record false positives, edited top-level artifacts, and added,
removed, or replaced stored references.  Edit counts describe deterministic
patches, not developer effort.

A separate \emph{contract-only} suite mutates typed relations absent from the
direct representation: it removes a capability-use provider, duplicates an
agent's source-agent identifier, or assigns a domain capability to the runtime
orchestrator.  Each mutation is applied once per case.  We report the canonical
V2 validator and the stricter pipeline agent--provider contract separately.
The baseline is marked ``not representable'' rather than failed; these results
measure added invariant scope, not comparative superiority on an impossible
task.

For descriptive cost, each adapter repeatedly constructs its normalized
responsibility view and enumerates all routes.  After four warm-ups, execution
order alternates for 40 measured repetitions per case using a monotonic
nanosecond clock.  The measurement includes adapter construction from already
parsed artifacts and route enumeration.  It excludes JSON I/O, v0.5-to-V2
generation, Unity, networking, and model latency.  We report medians,
quartiles, ranges, and the V2/direct median ratio without inferential claims.

\subsection{Reproducibility and analysis boundary}

The benchmark is API-free and records raw per-probe results, environment
metadata, source and input hashes, mutations, validation issues, timing samples,
and generated tables.  A verifier regenerates fresh V2 instances, reruns
targeted tests, checks the three Unity smokes when the Editor is available, and
scans the anonymous artifact for identifiers and secrets.

No confidence intervals or significance tests are attached to exhaustive,
deterministic probes over three purposively authored cases; doing so would imply
a sampling model the design does not have.  The evaluation does not measure
answer correctness, usefulness, trust, workload, preference, end-to-end
latency, or maintenance time.  It also does not use a language model as a
semantic judge.
